%% =====================================================================
%%  Corrigés — Unit 9 « Excuse me, where's the bank? »
%% =====================================================================
\corrigesUnite{9}{Excuse me, where's the bank?}

\begin{corrige}{1}
\begin{reponses}[3]
\item \rep{post office}
\item \rep{museum}
\item \rep{supermarket}
\item \rep{hotel}
\item \rep{station}
\item \rep{pharmacy}
\end{reponses}
\end{corrige}

\begin{corrige}{2}
\begin{reponses}[2]
\item \rep{Turn left.}
\item \rep{Go straight ahead.}
\item \rep{Don't take the bus.}
\item \rep{Take the subway.}
\item \rep{Cross the street.}
\item \rep{Don't be late!}
\end{reponses}
\pourquoi{l'impératif n'a pas de sujet ; le négatif se forme avec \en{Don't}, même pour \en{be} :
\en{Don't be late!}}
\end{corrige}

\begin{corrige}{3}
\begin{center}\otPlanVille[
  \otChemin{otGreenDark}{(73.5,2) -- (73.5,50) -- (62,50) -- (62,55)}{(62,58.5)}{1}
  \otChemin{otViolet}{(76,2) -- (76,22) -- (80,22)}{(84,22)}{2}
  \otChemin{otRedDark}{(77.5,2) -- (77.5,48.5) -- (90,48.5) -- (90,54)}{(94,57.5)}{3}
  \otChemin{otOrangeDark}{(74.5,2) -- (74.5,11) -- (60,11) -- (60,14)}{(56,17)}{4}]\end{center}
\begin{reponses}[4]
\item \rep{the park}
\item \rep{the supermarket}
\item \rep{the museum}
\item \rep{the post office}
\end{reponses}
\pourquoi{en marchant vers l'ouest (à gauche), le côté droit est le côté nord, et inversement. Suivez
chaque itinéraire avec le doigt sur le plan : les numéros correspondent aux phrases.}
\end{corrige}

\begin{corrige}{4}
\begin{reponses}[2]
\item \rep{next to}
\item \rep{opposite}
\item \rep{between}
\item \rep{on the corner of}
\item \rep{opposite}
\end{reponses}
\pourquoi{le musée et le parc sont séparés par \en{Second Avenue} : ils sont \en{opposite} (américain :
\en{across from}). L'hôtel touche à la fois \en{Third Avenue} et \en{Oak Street}.}
\end{corrige}

\begin{corrige}{5}
\begin{reponses}[6][french]
\item \rep{b}
\item \rep{d}
\item \rep{f}
\item \rep{a}
\item \rep{c}
\item \rep{e}
\end{reponses}
\end{corrige}

\begin{corrige}{6}
\begin{reponses}[1]
\item \rep{Excuse me, where's the station?}
\item \rep{Excuse me, how do I get to the museum?}
\item \rep{Excuse me, is there a bank near here?}
\item \rep{Is it far?}
\end{reponses}
\retiens{ajoutez \en{please} à la fin pour être encore plus poli : \en{Excuse me, where's the
station, please?}}
\end{corrige}

\begin{corrige}{7}
\begin{reponses}[4]
\item \rep{leave}
\item \rep{rich}
\item \rep{feet}
\item \rep{ship}
\item \rep{feel}
\item \rep{will}
\item \rep{least}
\item \rep{slip}
\end{reponses}
\pourquoi{\phon{iː} est nettement plus long et plus tendu que \phon{ɪ}. Texte de la piste
page~\pageref{trans-9.10}.}
\piege{prononcer \en{live} comme \en{leave} : « j'habite à Lyon » deviendrait « je quitte Lyon » !}
\end{corrige}

\begin{corrige}{8}
\begin{center}\otPlanVille[
  \otChemin{otNavy}{(74,2) -- (74,89) -- (122,89) -- (122,85)}{(126,82)}{1}
  \otChemin{otTeal}{(76,2) -- (76,9) -- (120,9) -- (120,14)}{(124,17)}{2}
  \otChemin{otGreenDark}{(75,2) -- (75,11) -- (26,11) -- (26,14)}{(22,17)}{3}]\end{center}
\begin{reponses}[3]
\item \rep{the hotel}
\item \rep{the pharmacy}
\item \rep{the bank}
\end{reponses}
\pourquoi{texte de la piste page~\pageref{trans-9.11}. Au 2, la pharmacie est « sur la gauche, à
côté du supermarché » quand on marche vers l'est sur \en{Main Street}.}
\end{corrige}

\begin{corrige}{9}
\begin{reponses}[4]
\item \rep{Excuse}
\item \rep{straight}
\item \rep{right}
\item \rep{far}
\item \rep{walk}
\item \rep{opposite}
\item \rep{don't}
\item \rep{Take}
\end{reponses}
\end{corrige}

\begin{corrige}{10}
\begin{reponses}[1]
\item \rep{Excuse me, where's the station, please?}
\item \rep{Go straight ahead, then turn right.}
\item \rep{Is it far? — No, it's a five-minute walk.}
\item \rep{The bank is next to the post office.}
\item \rep{Don't take the subway!}
\item \rep{Get off at Central Park.}
\item \rep{Is there a supermarket near here?}
\item \rep{The museum is on the left, opposite the park.}
\end{reponses}
\piege{phrase 2 : « à droite » = \en{right}, mais « tout droit » = \en{straight ahead}.}
\end{corrige}

\begin{corrige}{11}
Réponse modèle (vers la pharmacie) :
\begin{anglais}
\begin{tcolorbox}[otBase,colback=otGreenTint,colframe=otGreenTint,fontupper=\fontsize{10.4}{14}\selectfont]
Go straight ahead on Second Avenue to Main Street. Turn right on Main Street. Walk past the
supermarket. The pharmacy is on the left, on the corner of Third Avenue and Main Street.
\end{tcolorbox}
\end{anglais}
\retiens{\en{walk past} = « passer devant ». Pour votre propre itinéraire, réutilisez \en{go straight
ahead}, \en{turn left / right}, \en{it's on the left / right}, \en{next to}, \en{opposite}.}
\end{corrige}

\begin{corrige}{12}
Réponses modèles :
\begin{reponses}[1]
\item \rep{Go straight ahead and turn right on Oak Street. It's on the left.}
\item \rep{No, it isn't. It's a five-minute walk.}
\item \rep{Yes, there is. It's on Main Street, on the corner of Third Avenue.}
\item \rep{Don't take the bus. Walk!}
\item \rep{You're welcome!}
\end{reponses}
\end{corrige}

\begin{corrigetest}
\begin{reponses}[5][french]
\item \rep{a}
\item \rep{c}
\item \rep{b}
\item \rep{b}
\item \rep{a}
\item \rep{c}
\item \rep{a}
\item \rep{b}
\item \rep{c}
\item \rep{a}
\end{reponses}
\pourquoi{question 7 : dans \en{a five-minute walk}, \en{minute} reste au singulier, car l'expression
est un adjectif composé.}
\end{corrigetest}
